% MMW-METHOD-CONTRACT: 17e0e931dd01e167c81738be87e8eca8e00933987071026ea2be151827f2431e
% MMW-METHOD-CONTRACT: ead1dc31de5453e83924fd80c141e14521697c804489a275435dbb5338b47dac
% MMW-ID: CON-Q1-1
% MMW-ID: CON-Q1-2
% MMW-ID: CON-Q1-3
% MMW-ID: CON-Q1-4
% MMW-ID: CON-Q1-5
% MMW-ID: CON-Q1-6
% MMW-ID: CON-Q1-7
% MMW-ID: CON-Q2-1
% MMW-ID: CON-Q2-2
% MMW-ID: CON-Q2-3
% MMW-ID: CON-Q2-4
% MMW-ID: CON-Q2-5
% MMW-ID: CON-Q2-6
% MMW-ID: CON-Q3-1
% MMW-ID: CON-Q3-2
% MMW-ID: CON-Q3-3
% MMW-ID: CON-Q3-4
% MMW-ID: CON-Q3-5
% MMW-ID: CON-Q3-6
% MMW-ID: CON-Q3-7
% MMW-ID: CON-Q4-1
% MMW-ID: CON-Q4-2
% MMW-ID: CON-Q4-3
% MMW-ID: CON-Q4-4
% MMW-ID: CON-Q4-5
% MMW-ID: CON-Q4-6
% MMW-ID: OBJ-Q1
% MMW-ID: OBJ-Q2
% MMW-ID: OBJ-Q3
% MMW-ID: OBJ-Q4
% MMW-ALGORITHM: 严格闭环预测；五方案增量任务局部搜索；离散SOC动态规划；按情景动态重建邻域的任务-储能交替有限搜索
\section{模型的建立与求解}
\subsection{任务与能源统一约束}
令 $x_i=(r_i,s_i)$ 表示任务 $i$ 的执行区域与开工时刻，$f_i=s_i+d_i$。每个任务只执行一次；实时任务满足 $s_i=a_i$，弹性任务满足 $a_i\le s_i$ 且 $f_i\le\min(l_i,2406)$。区域间实际时延不得超过任务上限。对区域 $r$、小时 $t$，以区间重叠指示量累计 GPU 与瞬时 IT 功率，并要求
\[
G_{rt}\le \bar G_r,\qquad L^{\rm nonAI}_{rt}+L^{\rm AI}_{rt}\le \bar P^{\rm IT}_r,
\]
且乘以 PUE 后不超过设施功率边界。能源侧满足逐时守恒，购电、售电、充电、放电均非负，且购售电与充放电分别互斥。储能状态满足
\[
E_{r,t+1}=E_{rt}+\eta_r^c P^c_{rt}-P^d_{rt}/\eta_r^d,
\quad E_{r,2407}=E_{r0}.
\]
成本按购电支出减售电收入计算，碳排放只累计购电量与碳强度的乘积；负成本表示题面核算边界内售电收入高于购电支出，并不表示物理能耗为负。

\subsection{问题一：闭环预测}
对六区域三任务类型形成 18 条逐时 GPU 到达需求序列。候选包括日周期与周周期季节朴素模型，以及 $\epsilon\in\{1.10,1.35,1.75\}$、$\alpha\in\{10^{-4},10^{-3},10^{-2}\}$ 的 Huber 稳健模型\cite{huber1964robust}。训练、验证、测试边界分别为 2351、2352--2375 和 2376--2399；多步预测逐点回填预测值而非真实未来值。最终 17 条序列选择 Huber，1 条选择 168 小时季节朴素模型，平均 MAE=25.5892、RMSE=34.5070、WAPE=0.8231。图~\ref{fig:gpu} 展示最终任务方案的区域 GPU 利用率。
\begin{figure}[H]\centering
\includegraphics[width=.86\textwidth]{fig_1_gpu_utilization.png}
\caption{区域 GPU 利用率（完整任务轨迹复算）}\label{fig:gpu}
\end{figure}

\subsection{问题二：五方向有限邻域}
对成本、碳、新能源、服务和平衡五个方向分别检查 $10000$ 个任务移动候选。候选通过增量数组判断容量与能源变化，结束后再对 50000 个任务执行唯一性、时窗、时延、GPU、IT、设施功率和能源边界全量审计。表~\ref{tab:q2} 给出结果；方案名表示搜索方向，不表示跨方案或完整空间中的对应目标最优值。
\begin{table}[H]\centering\small
\caption{问题二五类搜索结果}\label{tab:q2}
\begin{tabular}{lrrrrr}\toprule
方向 & 成本/亿元 & 碳/t & 新能源率 & 时延/ms & 接受移动\\\midrule
成本 & -4.2334 & 10.2851 & 0.680386 & 6.2235 & 422\\
碳 & -4.1966 & 14.3916 & 0.679078 & 5.1594 & 0\\
新能源 & -4.2230 & 14.3916 & 0.680193 & 7.0758 & 761\\
服务 & -4.1966 & 14.3916 & 0.679078 & 5.1594 & 0\\
平衡 & -4.2180 & 16.5160 & 0.679962 & 6.5824 & 854\\\bottomrule
\end{tabular}
\end{table}
平衡方案迁移 561 个任务，服务质量为 0.990393。五次搜索合计恰好评价 50000 个候选，每个方向均单独通过一次全量任务审计。

\subsection{问题三：离散 SOC 动态规划与 Pareto 筛选}
每个区域使用 21 个等距 SOC 状态并显式加入初始 SOC，逐时枚举可达状态转移。候选包括 no-action、成本、碳、新能源、峰值、爬坡和平衡七类策略；相同轨迹仅计一次，并按多目标支配关系筛选\cite{miettinen1999nonlinear}。七个候选形成 6 条唯一轨迹，平衡方案通过候选集内非支配检验，成本、碳、新能源利用率、峰值和爬坡分别为 $-458584632.33$ 元、$0.224084$ tCO$_2$、0.771439、0.373986 MW 和 0.747972 MW。no-action 对应成本 $-419820327.74$ 元、碳排放 8.850563 tCO$_2$。图~\ref{fig:soc} 给出选中轨迹。
\begin{figure}[H]\centering
\includegraphics[width=.86\textwidth]{fig_2_storage_soc.png}
\caption{问题三选中储能轨迹的 SOC}\label{fig:soc}
\end{figure}

\subsection{问题四：动态任务--储能交替}
17 个情景包括基准、四个碳约束水平、三组峰谷价差、三组售电机制、三组新能源水平及种子 2026、2027、2028 的新能源波动。每个情景从相同任务基线和真实初始 SOC 独立重启；每轮先求储能，再把电价、碳强度、购电与弃电组合为任务移动信号，重新生成并评价 600 个任务候选。最多 8 轮，连续 2 轮无改善停止。

基准情景完成 2 轮、1200 次候选评价，成本 $-458341900$ 元、碳排放 0.578894 tCO$_2$、新能源利用率 0.768912、峰值购电 0.443765 MW；实时 SLA 和按期完成率均为 1，最大等待 6 h。10\% 碳削减情景在第 2 轮接受新任务方案，碳排放降至 0.424511 tCO$_2$；20\% 与 30\% 目标在声明启发式内未达到，不能据此判定数学不可行。全部情景合计 37 轮、22200 次候选评价，且每轮的候选与 incumbent 哈希均被记录。



\subsection{方法偏差与适用边界}
% MMW-LIMITATION: D1
\paragraph{方法边界 D1}Q2/Q4为预声明候选预算内最佳可行解，不是穷举最优解

% MMW-LIMITATION: D2
\paragraph{方法边界 D2}Q3 SOC使用有限状态离散，未计入题面未提供的退化成本和循环寿命硬上限

% MMW-LIMITATION: D3
\paragraph{方法边界 D3}Q4随机情景另使用固定种子2027和2028并逐情景独立重启

% MMW-LIMITATION: L1
\paragraph{方法边界 L1}无独立Oracle，最高可信等级为scenario-feasible

% MMW-LIMITATION: L2
\paragraph{方法边界 L2}储能SOC离散化且缺少退化成本/循环寿命数据

% MMW-LIMITATION: L3
\paragraph{方法边界 L3}Q2各命名方案表示搜索方向，不表示跨方案的对应目标最小值
